% !TEX root = ../UPDATED_main.tex
\begin{abstract}
Sequence saturation can erase the historical signal across a branch even when that branch has high conventional support. SatuTe addresses this problem directly by testing whether the subtrees separated by a branch remain statistically coherent. Using the original SatuTe test as the primary model, we first reproduce its Jukes--Cantor benchmark and show that bootstrap support can remain complete after branch-specific signal has fallen to the null rejection rate. We then evaluate two model-based extensions: an eigenvalue-weighted statistic that incorporates all persistent non-stationary directions of a reversible substitution model, and a corrected null for invariant-site mixtures. The weighted statistic reduces exactly to original SatuTe under Jukes--Cantor and improves detection mainly for short alignments near the long-branch decision boundary. Its largest matched GTR+F+G4 gain was 57.6 percentage points under taxon Bonferroni correction; across 1,728 matched protein cells, 831 differences were positive and none were negative. Independent truth-labeled experiments supported false-discovery control. Applications to Tree-of-Life and hepatitis C virus data localized broad losses of signal in 16S rRNA and codon-position-specific losses in viral coding regions. These results establish SatuTe as a complement to branch support and show how model-aware extensions can refine, without replacing, its branch-specific diagnosis of phylogenetic signal.
\end{abstract}
